GENRE:
Historiographic metafiction (the novel is metafictional in that it is about itself as a novel and historiographic in that it is about the tecgnoeconimic and sociocultural forces that drive history). Alternatively, though the novel contains both science fiction (low magic) and fantasy (high magic) it is a fabulation, defined by John Clute in the Encyclopedia of Science Fiction as 'any story which challenges the two main assumptions of genre sf: that the world can be seen; and that it can be told.'

TITLE:
The Visibilities

LOGLINE:
The travails of one man, and his avatars scattered across the multiverse, in attempting to overcome his cosmic loneliness and awake from the nightmare of history. 

COMP TITLES:
Henry Adams, 'A Law of Acceleration' amnd "The Rule of Phase Applied to History'
Kingsley Amis, Lucky Jim
Martin Amis, Money
Isaac Asimov, Extraterrestrial Civiisations
Margaret Atwood, The Handmaid's Tale
Julian Barnes, The History of the World in Ten and a Half Chapters
Jorge Luis Borges, Labyrinths
William Cooper, Scenes from Provinicial Life
Len Deighton, Bomber
BBC, Doctor Who
A.S. Byatt, Possession
John Clute and David Langford (eds), The Encylopedia of Science Fiction
Jonathan Coe, What a Carve Up!
Umberto Eco, The Name of the Rose
C.B.P. Finn, Thermal Physics
Michael Frayn, Sweet Dreams
Alan Garner, Red Shift
Garth Risk Hallberg, City on Fire
Douglas R. Hofstadter, Goedel, Escher, Bach: an Eternal Golden Braid
Uwe Johnson, Anniversaries
James Joyce, Ulysses
Donald E. Knuth, The TeXbook
Stanley Kubrick and Arthur C. Clarke, 2001: a Space Odyssey
Ursula K. Le Guin, A Wizard of Earthsea
David Lodge, How Far Can You Go? (US title: Souls and Bodies)
Mark W. Miller, Traveller Science-fiction Adventure in the Far Future
David Mitchell, Cloud Atlas
Christopher Priest, The Prestige
Thomas Pynchon, Gravity's Rainbow
Frank Richards, Billy Bunter and the Greyfriars Pretender
Olaf Stapledon, Star Maker
Laurence Sterne, The Life and Opinions of Tristram Shandy, Gentleman
Neal Stephenson, Cryptonomicon
Nicholas Tomalin and Ron Hall, The Strange Last Voyage of Donald Crowhurst
Jenny Turner, London Review of Books articles (https://www.lrb.co.uk/contributors/jenny-turner)
Vernor Vinge, A Fire Upon the Deep
Kurt Vonnegut Jr, Slaughterhouse-Five or the Children's Crusade: a Duty-dance with Death
David Foster Wallace, Infinite Jest
Evelyn Waugh, Sword of Honour trilogy
Wisden Cricketer's Almanack 1977, 114th edition

BACK COVER COPY:
The main text of The Visibilities is forced into a 366-element tensor that becomes so overloaded that inevitably it bursts out of its confines to infect the paratext of the novel including its peritext. The Visibilities is both a grimoire and a teratography. The front cover is a representation of a region of the Heptabrot set proceedureally generated with METAPOST. At the centre is the hexagonal 'fort' surrounded by beautifully rich, colourful, whirling, roiling sevenfold structure. The hexagonal fort is in the centre of the cover, the AR (Atomic Razor) device physically impressed (blind blocked) into the board (of the hardback, like the AP device on old Academic Press books). The back cover is similar, but shows a region of the Octabrot set. Here there is a eightfold structure surrounding the central, for this Multibrot set, heptagonal 'fort'. Ideally, matters will be arranfed so that that sevenfold and eightfold structures seamlessly meld into one another, but the realisation of that element of the book is left for the future. The front and spine will carry the name of the author and title (plus the publisher's device). There will no text on the back cover. The book will be sold new shrinkwrapped so that the barcode will go on the disposable outer covering so that it is not corrupted by any extraneous commercial clutter.     
